\chapter{G\'eom\'etrie alg\'ebrique et~g\'eom\'etrie~analytique}
\label{XII}

\begin{center}
{Mme M.\ \textsc{Raynaud}\footnote{D'apr\`es des notes in\'edites de
A.\, Grothendieck.}}
\end{center}
\vspace*{1cm}

Proc\'edant comme dans \cite{XII.10}, on associe \`a tout sch\'ema $X$
localement de type fini sur le corps des nombres complexes $\CC$ un
espace analytique~$X^\an$
\label{indnot:lb}dont l'ensemble sous-jacent est $X(\CC)$.

Dans les \No \ref{XII.2} et~\ref{XII.3} de cet expos\'e, nous donnons un
\og dictionnaire\fg entre les propri\'et\'es usuelles de~$X$ et de
$X^\an$ et entre les propri\'et\'es d'un morphisme $f\colon X\to
Y$ et du morphisme associ\'e $f^\an\colon X^\an \to Y^\an$.
\label{indnot:lc}
Nous montrons ensuite que les th\'eor\`emes de comparaison entre
faisceaux coh\'erents sur~$X$ et $X^\an$, \'etablis dans \cite[\No
12]{XII.10} pour une vari\'et\'e projective, sont encore valables lorsque
$X$ est un sch\'ema propre.

Enfin nous prouvons au \No \ref{XII.5} l'\'equivalence de la cat\'egorie des
rev\^etements \'etales finis de~$X$ et de la cat\'egorie des
rev\^etements \'etales finis de~$X^\an$. En prime au lecteur, nous
donnons une nouvelle d\'emonstration du th\'eor\`eme de
Grauert-Remmert \cite{XII.6}, utilisant la r\'esolution des
singularit\'es~\cite{XII.8}.

\section{Espace analytique associ\'e \`a un sch\'ema}
\label{XII.1}

Soit $X$ un sch\'ema localement de type fini sur~$\CC$. Soit $\Phi$
le foncteur de la cat\'egorie des espaces analytiques \cite[\No 9]{XII.4} dans
la cat\'egorie des ensembles, qui \`a un espace analytique~$\othercal{X}$ associe l'ensemble des morphismes d'espaces annel\'es en
$\CC$-alg\`ebres $\Hom_\CC (\othercal{X},X)$. On a le th\'eor\`eme
suivant:

\begin{theoremedefinition}
\label{XII.1.1} Le foncteur $\Phi$ est repr\'esentable par un
espace analytique $X^\an$ et un morphisme $\varphi\colon X^\an\to
X$. On dit que $X^\an$ est l'espace analytique associ\'e \`a~$X$.
\index{espace analytique associ\'e|hyperpage}
Si $|X^\an|$ est l'ensemble sous-jacent \`a $X^\an$, $\varphi$
induit une bijection de~$|X^\an|$ sur l'ensemble $X(\CC)$ des points
de~$X$ \`a valeurs dans $\CC$. De plus, pour chaque point $x$ de~$X^\an$, le morphisme
$$
\varphi_x\colon\cal{O}_{X,\varphi(x)}\to\cal{O}_{X^\an,x},
$$
qui est n\'ecessairement local, donne par passage aux
compl\'et\'es un isomorphisme
$$
\widehat{\varphi}_x\colon\widehat{\cal{O}}_{X,\varphi(x)}\isomto
\widehat{\cal{O}}_{X^\an,x},
$$
En particulier le morphisme $\varphi$ est plat.
\end{theoremedefinition}

Notons que le fait que $\varphi$ induise une bijection de~$X^\an$ sur~$X(\CC)$ r\'esulte de la propri\'et\'e universelle de~$X^\an$.
D'autre part on a les assertions suivantes:

\begin{enumerate}
\item [a)] Si le th\'eor\`eme est vrai pour un sch\'ema $Y$, il
en est de m\^eme pour tout sous-sch\'ema $X$ de~$Y$. Supposons
d'abord que $X$ soit un sous-sch\'ema ouvert de~$Y$; si $\psi\colon
Y^\an\to Y$ est le morphisme canonique, $\psi^{-1}(X)$ est un ouvert
de~$Y^\an$ que l'on muni de la structure d'espace analytique induite
par celle de~$Y^\an$. Comme tout morphisme d'un espace analytique
$\othercal{X}$ dans $X$ se factorise \`a travers $Y^\an$ d'apr\`es la
propri\'et\'e universelle de ce dernier, donc \`a travers
$X^\an$ qui est le produit fibr\'e $Y^\an\times_Y X$, $X^\an$ est
l'espace analytique associ\'e \`a $X$. Enfin l'assertion
concernant les $\varphi_x$ est \'evidente.

Il suffit maintenant de consid\'erer le cas o\`u $X$ est un
sous-sch\'ema ferm\'e de~$Y$. Soit~$I$ le $\cal{O}_Y$-Id\'eal
coh\'erent d\'efinissant $X$; alors $I\cdot\cal{O}_{Y^\an}$ est un
faisceau coh\'erent d'id\'eaux sur~$\cal{O}_{Y^\an}$ qui
d\'efinit un sous-espace analytique ferm\'e $X^\an$ de~$Y^\an$; on
voit comme dans le cas d'un sous-sch\'ema ouvert que $X^\an$ est
l'espace analytique associ\'e \`a $X$. Soit $\varphi\colon
X^\an\to X$ le morphisme canonique. Pour tout point $x$ de~$X^\an$, le
morphisme $\varphi_x$ n'est autre que le morphisme
$$
\cal{O}_{Y,\psi(x)}/I_{\psi(x)}\to
\cal{O}_{Y^\an,x}/I_{\psi(x)}\cdot\cal{O}_{Y^\an,x}
$$
induit par $\psi_x$; son compl\'et\'e
$$
\widehat{\varphi}_x\colon
\widehat{\cal{O}}_{Y,\psi(x)}/I_{\psi(x)}\cdot\widehat{\cal{O}}_{Y,\psi(x)}
\to\widehat{\cal{O}}_{Y^\an,x}/I_{\psi(x)}\cdot\widehat{\cal{O}}_{Y^\an,x}
$$
est un isomorphisme puisque $\widehat\psi_x$ en est un, ce qui
d\'emontre a).

\item [b)] Si l'on a deux $\CC$-sch\'emas $X_1$, $X_2$, tels que
$X_1^\an$ et $X_2^\an$ existent, alors il en est de m\^eme de
$(X_1\times X_2)^\an$. Soient en effet $\varphi_1\colon X_1^\an\to
X_1$, $\varphi_2\colon X_2^\an\to X_2$ les morphismes canoniques,
$p_1$, $p_2$ les deux projections de~$X_1^\an\times X_2^\an$. On
d\'eduit formellement de EGA~I~1.8.1 que $X_1\times X_2$ est le
produit de~$X_1$ et $X_2$ dans la cat\'egorie des espaces
annel\'es en anneaux locaux; il en r\'esulte que les morphismes
$\varphi_1\cdot p_1$ et $\varphi_2\cdot p_2$ d\'efinissent un
morphisme $\varphi\colon X_1^\an\times X_2^\an\to X_1\times X_2$ et
que le couple $(X_1^\an\times X_2^\an,\varphi)$ repr\'esente le
foncteur $\othercal{X}\mto\Hom_\CC(\othercal{X},X_1\times X_2)$.

\item [c)] Si $\othercal{E}^1$ d\'esigne l'espace affine de dimension 1,
\ie l'espace topologique $\CC$ muni du faisceau des fonctions
holomorphes, le foncteur{\renewcommand{\thefootnote}{*}\addtocounter{footnote}{-1}\footnote{\lcrochetbf Ajout\'e en 2003: $E^1_\CC$ d\'esigne la droite affine (alg\'ebrique) sur $\CC$.\rcrochetbf}}
$\othercal{X}\mto\Hom_\CC (\othercal{X},E^1_\CC)$
est repr\'esentable par $\othercal{E}^1$, le morphisme canonique
$\varphi\colon\othercal{E}^1\to E^1_\CC$ \'etant le morphisme
\'evident. En effet se donner un morphisme d'un espace analytique
$\othercal{X}$ dans $E^1_\CC$ \'equivaut \`a se donner un
\'el\'ement de~$\Gamma(\othercal{X},\cal{O}_\othercal{X})$, ce qui revient
aussi \`a se donner un morphisme de~$\othercal{X}$ dans~$\othercal{E}^1$. On
a \'evidemment une bijection $|\othercal{E}^1|\simeq E^1(\CC)$, et, pour
chaque point $x\in\othercal{E}^1$, le morphisme $\widehat\varphi_x$ n'est
autre que le morphisme identique d'un anneau de s\'eries formelles
\`a une variable sur~$\CC$.
\end{enumerate}

\enlargethispage{.5cm}On d\'eduit de b) et c) que le th\'eor\`eme est vrai pour
l'espace affine $E^n_\CC$, $n\geq 0$. Utilisant a), on voit qu'il en
est de m\^eme pour tout sch\'ema affine $X$, localement de type
fini sur~$\CC$. Si l'on ne suppose plus $X$ affine et si $(X_i)$ est
un recouvrement de~$X$ par des ouverts affines, il r\'esulte de la
propri\'et\'e universelle et de a) que les $X_i^\an$ se recollent
et d\'efinissent ainsi l'espace analytique $X^\an$ associ\'e
\`a~$X$.

\subsection{}
\label{XII.1.2}
Soit $f\colon X\to Y$ un morphisme de~$\CC$-sch\'emas localement de
type fini. Si $\varphi\colon X^{\an}\to X$ et $\psi\colon Y^{\an}\to
Y$ sont les morphismes canoniques, il r\'esulte de la
propri\'et\'e universelle de~$Y^{\an}$ qu'il existe un unique
morphisme $f^{\an}\colon X^{\an}\to Y^{\an}$ tel que le diagramme
$$
\xymatrix{{X^{\an}}\ar[r]\ar[d]_-{f^{\an}}&{X}\ar[d]^-f\\{Y^{\an}}\ar[r]&{Y}}
$$
soit
commutatif. On a donc d\'efini un foncteur $\Phi$ de la
cat\'egorie des $\CC$-sch\'emas localement de type fini dans la
cat\'egorie des espaces analytiques.

Le foncteur $\Phi$ commute aux limites projectives finies. Il suffit
en effet de voir que $\Phi$ commute aux produits fibr\'es. Or, si
$X$, $Y$, $Z$, sont des sch\'emas localement de type fini sur~$\CC$,
il r\'esulte du fait que $X\times_Z Y$ est le produit fibr\'e de
$X$ et $Y$ au-dessus de~$Z$ dans la cat\'egorie des espaces
annel\'es en anneaux locaux que $X^{\an}\times_{Z^{\an}}Y^{\an}$
satisfait \`a la propri\'et\'e universelle qui caract\'erise
$(X\times_Z Y)^{\an}$.

\subsection{}
\label{XII.1.3}
Soient $X$ un $\CC$-sch\'ema localement de type fini, $X^{\an}$
l'espace analytique associ\'e, $\varphi\colon X^{\an}\to X$ le
morphisme canonique. Si $F$ est un ${\cal{O}}_X$-Module, l'image
inverse $\varphi^*F=F^{\an}$ est un faisceau de modules sur~${\cal{O}}_{X^{\an}}$. On d\'efinit ainsi un foncteur de la
cat\'egorie des $\cal{O}_X$-Modules dans la cat\'egorie des
Modules sur~$X^{\an}$. Ce foncteur commute aux limites inductives
(EGA~0~4.3.2). Le faisceau $\cal{O}_{X^{\an}}$ \'etant coh\'erent
\cite[\no18~\S 2~th.\, 2]{XII.4}, il transforme faisceaux coh\'erents en
faisceaux coh\'erents (EGA~0~5.3.11). On a de plus:

\begin{subproposition}
\label{XII.1.3.1}
Le foncteur qui \`a un $\cal{O}_{X^{\an}}$-Module $F$ associe son
image inverse $F^{\an}$ sur~$X^{\an}$ est exact, fid\`ele,
conservatif.
\end{subproposition}

L'exactitude r\'esulte du fait que le morphisme $\varphi\colon
X^{\an}\to X$ est plat \eqref{XIII.1.1}. Prouvons que le foncteur $F\mto
F^{\an}$ est fid\`ele. Compte tenu de l'exactitude, il suffit de
montrer que, si $F^{\an}$ est nul, il en est de m\^eme de~$F$. Or,
pour tout point $x$ de~$X^{\an}$, on a alors $F_{\varphi
(x)}\otimes_{\cal{O}_{X,\varphi (x)}}\cal{O}_{X^{\an},x}=0$. Le
morphisme $\cal{O}_{X,\varphi (x)}\to \cal{O}_{X^{\an},x}$ \'etant
fid\`element plat, on a $F_{\varphi (x)}=0$ pour tout point
ferm\'e $\varphi (x)$ de~$X$, et, comme $X$ est de Jacobson
(EGA~IV~10.4.8), ceci implique que $F$ est nul.

Le
fait que le foncteur $F\mto F^{\an}$ soit conservatif est formel
\`a partir de l'exactitude et de la fid\'elit\'e.

\section[Comparaison des propri\'et\'es d'un sch\'ema]{Comparaison des propri\'et\'es d'un sch\'ema et de l'espace
analytique \hbox{associé}}
\label{XII.2}

\begin{proposition}
\label{XII.2.1}
Soient $X$ un $\CC$-sch\'ema localement de type fini, $X^{\an}$
l'espace analytique associ\'e, $n$ un entier. Consid\'erons la
propri\'et\'e $P$ d'\^etre
\begin{itemize}
\item[(i)] non vide
\item[(i')] discret
\item[(ii)] de Cohen-Macaulay
\item[(iii)] $(\mathrm{S}_n)$
\item[(iv)] r\'egulier
\item[(v)] $(\mathrm{R}_n)$
\item[(vi)] normal
\item[(vii)] r\'eduit
\item[(viii)] de dimension $n$.
\end{itemize}
Alors, pour que $X$ poss\`ede la propri\'et\'e $P$, il faut et
il suffit qu'il en soit ainsi de~$X^{\an}$.
\end{proposition}

Soit $\varphi\colon X^{\an}\to X$ le morphisme
canonique. (i)~r\'esulte du fait que l'on a $|X^{\an}|=X(\CC )$
\eqref{XIII.1.1} et du fait que $X$ est de Jacobson (EGA~IV~10.4.8). Dire que~$X$
(\resp $X^{\an}$) est discret \'equivaut \`a dire que l'on a $\dim
X=0$ (\resp $\dim X^{\an}=0$ d'apr\`es \cite[\no19~\S 4~cor.\, 6]{XII.4}); (i')
r\'esulte donc de~(viii).

Soit $P$ l'une des propri\'et\'es (ii) \`a (vii). Pour que $X$
poss\`ede la propri\'et\'e $P$, il faut et il suffit que $P$
soit v\'erifi\'ee en chaque point ferm\'e de~$X$; en effet, $X$
\'etant excellent (EGA~IV~7.8.6~(iii)), l'ensemble des points
o\`u $X$ v\'erifie $P$ est un ouvert (\loccit) et, si cet ouvert
contient tous les points ferm\'es, il est \'egal \`a $X$ tout
entier. Dire que $X$ (\resp $X^{\an}$) a la propri\'et\'e $P$
\'equivaut donc \`a dire que, pour tout point $x$ de~$X^{\an}$,
l'anneau local $\cal{O}_{X,\varphi (x)}$ (\resp $\cal{O}_{X^{\an},x}$)
a la propri\'et\'e $P$. Comme le fait qu'un anneau local excellent
ait la propri\'et\'e $P$ se voit apr\`es passage au
compl\'et\'e, la proposition r\'esulte des isomorphismes
$\widehat{\cal{O}}_{X,\varphi(x)}\isomto\widehat{\cal{O}}_{X^{\an},x}$
dans les cas (ii) \`a (vii). Il en est de m\^eme dans le cas
(viii), compte tenu des relations
$$
\dim X=\sup_x~\dim \cal{O}_{X,\varphi (x)}\qquad \dim
X^{\an}=\sup_x~\dim \cal{O}_{X^{\an},x}
$$
o\`u $x\in X^{\an}$. Ceci ach\`eve la d\'emonstration.

\begin{proposition}
\label{XII.2.2}
Soient $X$ un $\CC$-sch\'ema localement de type fini,
$\varphi\colon X^{\an}\to X$ le morphisme canonique, $T$ une partie
localement constructible de~$X$. Alors on a la relation
$$
\varphi^{-1}(\overline{T})=\overline{\varphi^{-1}(T)}.
$$
\end{proposition}

On peut supposer que $T$ est un ouvert dense de~$X$. Soit $H$ le
sous-sch\'ema ferm\'e r\'eduit de~$X$ d'espace sous-jacent
$X-T$; l'espace associ\'e $H^{\an}$ est un sous-espace analytique
ferm\'e de~$X^{\an}$ d'espace sous-jacent
$X^{\an}-\varphi^{-1}(T)$. On doit montrer que tout point $x$ de
$H^{\an}$ appartient \`a $\overline{\varphi^{-1}(T)}$. Or, en un tel
point $x$, le germe d'espace analytique $(X^{\an},x)$ contient le
sous-germe $(H^{\an},x)$, et celui-ci est d\'efini par un Id\'eal
non nilpotent de~$\cal{O}_{X^{\an},x}$. Il r\'esulte alors du
Nullstellensatz \cite[\no19~\S 4~cor.\, 3]{XII.4} que tout voisinage ouvert de~$x$
contient des points de~$X^{\an}$ qui n'appartiennent pas \`a
$H^{\an}$, ce qui prouve bien que l'on a
$x\in\overline{\varphi^{-1}(T)}$.

\begin{corollaire}
\label{XII.2.3}
Soient
$X$ un $\CC$-sch\'ema localement de type fini, $\varphi
\colon X^\an \to X$ le morphisme canonique, $T$ une partie localement
constructible de~$X$. Pour que $T$ soit une partie ouverte (\resp une
partie ferm\'ee, \resp une partie dense), il faut et il suffit qu'il
en soit ainsi de~$\varphi^{-1} (T)$.
\end{corollaire}

Le corollaire r\'esulte de~\ref{XII.2.2} et du fait que, $X$
\'etant un sch\'ema de Jacobson (EGA~IV~10.4.8), deux parties
localement constructibles de~$X$ qui ont m\^eme trace sur l'ensemble
tr\`es dense $X (\CC )$ sont \'egales.

\begin{proposition}
\label{XII.2.4} Soit $X$ un $\CC$-sch\'ema localement de type
fini. Pour que $X$ soit connexe (\resp irr\'eductible), il faut et
il suffit qu'il en soit ainsi de~$X^\an$.
\end{proposition}

Supposons $X^\an $ connexe (\resp irr\'eductible). L'image $X (\CC)
$ de~$X^\an$ dans $X$ est alors connexe (\resp irr\'eductible). Il
en r\'esulte que $X$ est connexe (\resp irr\'eductible) car les
parties ferm\'ees de~$X$ et $X (\CC)$ se correspondent bijectivement
(EGA~IV~10.1.2).

Inversement supposons $X$ connexe (\resp irr\'eductible), et
montrons qu'il en est de m\^eme de~$X^\an$. On peut se borner au cas
o\`u $X$ est irr\'eductible. Supposons en effet $X$ connexe.
\'Etant donn\'e un point $x$ de~$X$, l'ensemble des points $y \in X$
tels qu'il existe une suite finie de sous-sch\'emas ferm\'es
irr\'eductibles $X_1, \dots, X_n$ de~$X$, avec $x \in X_1$, $y \in
X_n$, $X_i \cap X_{i+1} \neq \emptyset $ pour $1 \leq i \leq n-1$,
est un ensemble \`a la fois ouvert et ferm\'e, donc \'egal \`a
$X$ tout entier. Pour une suite $X_1, \dots, X_n$ telle que
pr\'ec\'edemment, on a aussi $X_{i}^\an \cap X_{i+1}^\an \neq
\emptyset $ pour $1 \leq i \leq n-1 $; si l'on suppose
d\'emontr\'e que les $X_{i}^\an$ sont connexes, il en est alors de
m\^eme de~$X^\an$.

On
suppose d\'esormais $X$ irr\'eductible. On peut supposer de plus
$X$ affine. En effet, si $(U_i)_{i \in I}$ est un recouvrement de~$X$
par des ouverts affines, deux de ces ouverts ont une intersection non
vide, et la m\^eme propri\'et\'e est donc vraie pour le
recouvrement $(U_{i}^\an)_{i \in I} $ de~$X^\an$; si l'on suppose
d\'emontr\'e que les $U_{i}^\an $ sont irr\'eductibles, il en
est alors de m\^eme de~$X^\an$.

On peut supposer de plus que $X$ est normal. Soit en effet $\tilde X$
le normalis\'e de~$X$; comme le morphisme $\tilde X \to X$ est
surjectif, il est de m\^eme de~$\tilde X^\an \to X^\an$, ce qui
prouve que, si $\tilde X^\an $ est irr\'eductible, il est de
m\^eme de~$X^\an$.

On suppose d\'esormais $X$ affine normal. Comme les anneaux locaux
de~$X^\an$ sont int\`egres, il revient au m\^eme de dire que
$X^\an$ est irr\'eductible ou qu'il est connexe. En effet, si
$\cal{F}$ est une partie analytique ferm\'ee de~$X^\an$, l'ensemble
des points $x$ de~$X^\an$ o\`u l'on a $\codim_x (\cal{F}, X^\an ) =
0$ est un sous-ensemble analytique ferm\'e de~$X^\an$ \cite[\no 20 A
Cor.\, 1]{XII.4} qui est aussi ouvert; si $X^\an$ est connexe, ceci prouve que,
si l'on a $\cal{F} \neq X^\an$, $\cal{F}$ est rare, donc que $X^\an$
est irr\'eductible. On est ainsi ramen\'e \`a montrer que
$X^\an$ est connexe.

Soit
$$
i \colon X \to P
$$
une compactification de~$X$, o\`u $P$ est un $\CC$-sch\'ema
projectif normal et $i$ une immersion ouverte dominante. Il
r\'esulte alors de \cite[\no 12 th.\, 1]{XII.10} que $P^\an$ est connexe. Comme
$X^\an$ est obtenu en enlevant \`a $P^\an$ une partie analytique
ferm\'ee rare, il r\'esulte de~\ref{XII.2.5} ci-dessous que $X^\an$
est aussi connexe.

\begin{lemme}
\label{XII.2.5}
Soient
$\cal{P} $ un espace analytique normal connexe, $\cal{Y}$ une
partie analytique ferm\'ee rare, alors $\cal{X} = \cal{P} - \cal{Y}$
est connexe.
\end{lemme}

Lorsque $\cal{Y}$ est de codimension $\geq 2$, la proposition
r\'esulte de \cite[\no 3 prop.\, 4]{XII.11}. Dans le cas g\'en\'eral on peut
supposer, quitte \`a enlever \`a $\cal{P}$ une partie analytique
ferm\'ee de codimension $\geq 2$, que $\cal{P}$ et $\cal{Y}$
(consid\'er\'e comme sous-espace analytique r\'eduit de~$\cal{P}$) sont r\'eguliers. D'apr\`es le th\'eor\`eme des
fonctions implicites, tout point $y$ de~$\cal{Y}$ poss\`ede un
voisinage $\cal{U}$ isomorphe \`a une boule d'un espace affine
$\cal{E}^{n}$, de sorte que $\cal{U} \cap \cal{Y}$ soit d\'efini par
l'annulation d'un certain nombre de fonctions coordonn\'ees. Ceci
prouve que $\cal{U} - \cal{U} \cap \cal{Y}$ est connexe, et il en est
donc de m\^eme de~$\cal{X}$.

\begin{corollaire}
\label{XII.2.6} Soit $X$ un $\CC$-sch\'ema localement de type
fini; le morphisme
$$
\pi_0 (X^{\an}) \to \pi_0 (X)
$$
induit par le morphisme canonique $X^\an \to X$ est bijectif.
\end{corollaire}

\section{Comparaison des propri\'et\'es des morphismes}
\label{XII.3}

\begin{proposition}
\label{XII.3.1} Soient $f \colon X \to Y$ un morphisme de
$\CC$-sch\'emas localement de type fini, $f^\an \colon X^\an \to
Y^\an$ le morphisme d\'eduit de~$f$ sur les espaces analytiques
associ\'es. Soit $P$ la propri\'et\'e d'\^etre
\begin{enumerate}
\item[(i)] plat
\item[(ii)] net (\ie non ramifi\'e)
\item[(iii)] \'etale
\item[(iv)] lisse
\item[(v)] normal
\item[(vi)]
r\'eduit
\item[(vii)] injectif
\item[(viii)] s\'epar\'e
\item[(ix)] un isomorphisme
\item[(x)] un monomorphisme
\item[(xi)] une immersion ouverte.
\end{enumerate}
Alors, pour que $f$ poss\`ede la propri\'et\'e $P$, il faut et
il suffit qu'il en soit ainsi de~$f^\an$.
\end{proposition}

Notons $\varphi \colon X^\an \to X $ et $\psi \colon Y^\an \to Y$ les
morphismes canoniques. Soient $x$ un point de~$X^\an$, $y = f^\an
(x)$. Les morphismes $\cal{O}_{Y^\an, y} \to {\mathcal O}_{X^\an, x}$
et $\cal{O}_{Y, \psi(y)} \to \cal{O}_{X, \varphi(x)} $ d\'eduits de
$f$ et $f^{\an}$ donnent le m\^eme morphisme par passage aux
compl\'et\'es \eqref{XII.1.1}. D'apr\`es \cite[ch.\, 3 \S 5 prop.\, 4]{XII.2}
(\resp EGA~IV 17.4.4) il revient donc au m\^eme de dire que $f^\an $
v\'erifie la propri\'et\'e (i) (\resp (ii)) ou de dire que $f$
v\'erifie (i) (\resp (ii)) en chaque point ferm\'e de~$X$. Comme
l'ensemble des points de~$X$ o\`u (i) (\resp (ii)) est
v\'erifi\'e est un ouvert (EGA~IV 11.1.1 et I 3.3), ceci
d\'emontre (i) et (ii), donc aussi (iii).

\enlargethispage{.5cm}Soit $P$ la propri\'et\'e (iv) (\resp (v), \resp (vi)). Compte tenu
de~\ref{XII.2.1} ((v), (vi), (vii)), il revient au m\^eme de dire que
les fibres g\'eom\'etriques de~$f^\an$ aux diff\'erents points
$y$ de~$Y^\an$ sont r\'eguli\`eres (\resp normales, \resp r\'eduites) ou qu'il en est ainsi des fibres g\'eom\'etriques de
$f$ aux diff\'erents points ferm\'es $\psi (y)$ de~$Y$. Les cas
(iv) (\resp (v), \resp (vi)) r\'esultent alors de (i) et du fait que
l'ensemble des points de~$Y$ o\`u les fibres g\'eom\'etriques de
$f$ sont r\'eguli\`eres est un ouvert (EGA~IV 12.1.7).

(vii). Si $f$ est injectif, il en est de m\^eme de~$f^\an$.
Inversement supposons $f^\an$ injectif et montrons qu'il en est de
m\^eme de~$f$. On peut supposer
$f$ de type fini. Le morphisme $f^\an$ \'etant injectif, les fibres
de~$f$ aux points ferm\'es de~$Y$ sont radicielles; comme l'ensemble
des points de~$Y$ dont la fibre est radicielle est localement
constructible (EGA~IV 9.6.1) et comme $Y$ est un sch\'ema de
Jacobson, $f$ a toutes ses fibres radicielles donc est injectif.

(viii). Soient $\Delta \colon X \to X \times_Y X $ et $\Delta^\an
\colon X^\an \to X^\an \times_{Y^\an} X^\an $ les immersions
diagonales, $\Theta \colon X^\an \times_{Y^\an} X^\an \to X \times_Y X
$ le morphisme canonique. En vertu de~\ref{XII.2.3} il revient au
m\^eme de dire que $\Delta (X)$ est ferm\'e dans $X \times_Y X $
ou que $\Delta^\an (X^\an)$ est ferm\'e dans $X^\an \times_{Y^\an}
X^\an$.

Comme une immersion ouverte n'est autre qu'un morphisme \'etale
injectif (EGA~IV 17.9.1 et \cite[\no 13 \S 1]{XII.4}), (xi) r\'esulte de (iii)
et de (vii). Un isomorphisme \'etant la m\^eme chose qu'une
immersion ouverte surjective, (ix) r\'esulte de (xi) et de
(\ref{XII.3.2}~(i)) ci-dessous. Dire que $f$ est un monomorphisme
\'equivaut \`a dire que la morphisme diagonal $\Delta \colon X \to
X \times_Y X$ est un isomorphisme, donc (x) r\'esulte de (ix).

\begin{proposition}
\label{XII.3.2}
Soient $X$ et $Y$ deux $\CC$-sch\'emas localement de type fini,
$f\colon X \to Y$ un morphisme de type fini, $f^\an \colon X^\an \to
Y^\an $ le morphisme d\'eduit de~$f$ sur les espaces analytiques
associ\'es. Soit $P$ la propri\'et\'e d'\^etre
\begin{enumerate}
\item[(i)] surjectif
\item[(ii)] dominant
\item[(iii)] une immersion ferm\'ee
\item[(iv)] une immersion
\item[(v)] propre\footnote{Nous dirons qu'un morphisme d'espaces
analytiques est propre s'il l'est au sens de \cite[ch.\, 1 \S 10 \no 1]{XII.1} et
s'il est s\'epar\'e.}
\item[(vi)]
fini.
\end{enumerate}
Alors, pour que $f$ poss\`ede la propri\'et\'e $P$, il faut et
il suffit qu'il en soit ainsi de~$f^\an$.
\end{proposition}

Soient $\varphi \colon X^\an \to X$ et $\psi \colon Y^\an \to Y$ les
morphismes canoniques.

(i). Si $f$ est surjectif, pour tout point $y$ de~$Y^\an$, $f^{-1}
(\psi (y)) $ est une partie ferm\'ee non vide de~$X$; elle contient
donc au moins un point ferm\'e, ce qui prouve que $f^\an$ est
surjectif. Inversement, si $f^\an$ est surjectif, $f(X)$ est une
partie localement constructible de~$Y$ (EGA~IV 1.8.4) qui contient
tous les points ferm\'es de~$Y$; on a donc $f(X)=Y$.

(ii) r\'esulte de~\ref{XII.2.2}.

(iii). Si $f$ est une immersion ferm\'ee, il en est de m\^eme de
$f^\an$ d'apr\`es \ref{XII.1.1}~a). Inversement, si $f^\an$ est une
immersion ferm\'ee, il en est de m\^eme de~$f$ d'apr\`es
\ref{XII.3.1}~(x) et \ref{XII.3.2}~(v), car cela revient \`a dire
que $f$ est un monomorphisme propre (EGA~IV 8.11.5).

(iv). Il est clair que, si $f$ est une immersion, il en est de
m\^eme de~$f^\an$. Inversement supposons que $f^\an$ soit une
immersion, et soient $T$ l'image de~$X$ dans $Y$, $\overline{T}$
l'adh\'erence sch\'ematique de~$f$. On a une factorisation de~$f$,
$$
X \lto{i} \overline{T} \lto{j} J \quoi,
$$
o\`u $j$ est une immersion ferm\'ee, $i$ le morphisme canonique,
et on en d\'eduit la factorisation suivante de~$f^\an$
$$
X^\an \lto{i^\an} \overline{T}^{ \an} \lto{j^\an} Y^\an \
.
$$
Comme
$T = f(X)$ est une partie localement constructible de~$Y$ (EGA~IV
1.8.4), on~a, d'apr\`es~\ref{XII.2.2}, $\overline{T}^{ \an} =
\overline{f^\an (X^\an )}$. Il en r\'esulte que $i^\an (X^\an )$ est
un ouvert de~$\overline{T}^{ \an}$, donc que $i(X) $ est un ouvert de
$\overline{T}$. On consid\`ere la factorisation canonique de~$i$
$$
X \lto{i_1} i (X) \lto{i_2} \overline{T}.
$$
Le morphisme $i_1^\an $ est un monomorphisme propre, donc il en est de
m\^eme de~$i_1$ d'apr\`es \ref{XII.3.2}~(v) et \ref{XII.3.1}~(x);
ceci prouve que $i_1$ donc aussi $f$ est une immersion.

(v). Supposons que $f$ soit propre et montrons qu'il en est de
m\^eme de~$f^\an$. Le fait que $f^\an $ soit propre \'etant local
sur~$Y^\an$, on peut supposer $Y$ affine. D'apr\`es le lemme de Chow
(EGA~II 5.6.1), on peut trouver un $Y$-sch\'ema projectif $X'$ et un
morphisme projectif surjectif
$$
g \colon X' \to X.
$$
Le morphisme $(fg)^\an = f^\an g^\an $ est projectif donc propre,
$g^\an $ est surjectif, et il r\'esulte de \cite[ch.\, 1 \S 10]{XII.1} que
$f^\an $ est propre.

Inversement supposons $f^\an $ propre et montrons qu'il en est de
m\^eme de~$f$. D'apr\`es \ref{XII.3.1}~(viii) $f$ est
s\'epar\'e. Il reste \`a prouver que $f$ est universellement
ferm\'e, et il suffit m\^eme de montrer que $f$ est ferm\'e; en
effet, pour tout $Y$-sch\'ema $Y'$ localement de type fini, le
morphisme
$$
f_{(Y')} = h \colon X \times_Y Y' \to Y'
$$
sera aussi ferm\'e puisque $h^\an $ est propre. Soit $T$ une partie
ferm\'ee de~$X$; $f(T)$ est un ensemble localement constructible, et
l'on a
$$
f^\an (\varphi^{-1} (T)) = \psi^{-1} (f(T)).
$$
Comme
$f^\an $ est propre, $\psi^{-1} (f (T))$ est une partie ferm\'ee de
$Y^\an $, et il r\'esulte donc de~\ref{XII.2.2} que l'on a
$$
\psi^{-1} (\overline{f(T)}) = \psi^{-1} (f(T))\quoi .
$$
Cela entra\^ine que l'on a $f(T) = \overline{f(T)}$, \ie que $f$
est ferm\'ee donc que $f$ est propre.

(vi). Il revient au m\^eme de dire qu'un morphisme est fini ou qu'il
est propre \`a fibres finies (EGA~III 4.4.2 et \cite[\no 19 \S 5]{XII.4}).
Comme l'ensemble des points o\`u les fibres de~$f$ sont finies est
localement constructible (EGA~IV 9.7.9), les fibres de~$f$ sont finies
si et seulement si il en est ainsi des fibres de~$f^\an $; (vi)
r\'esulte donc de (v).

\begin{remarque}
\label{XII.3.3}
\begin{enumerate}
\item[a)] soit $f \colon X \to Y$ un morphisme de~$\CC$-sch\'emas
localement de type fini. Le fait que $f^\an $ soit un isomorphisme
local n'entra\^ine pas qu'il en soit de m\^eme de~$f$. En effet,
si $f$ est \'etale, $f^\an $ est \'etale donc est un isomorphisme
local \hbox{\cite[\no 13 \S 1]{XII.4}}, mais il n'en est pas n\'ecessairement ainsi
de~$f$.
\item[b)]
L'\'enonc\'e~\ref{XII.3.2}
n'est pas vrai si l'on ne suppose
pas $f$ de type fini. Montrons par exemple que $f^\an $ peut \^etre
une immersion ferm\'ee sans qu'il en soit de m\^eme de~$f$. Il
suffit en effet de prendre pour $X$ la somme de~$\ZZ $ copies de
$\Spec \CC$, et pour $Y$ la droite affine, et pour $f$ le morphisme
obtenu en envoyant les points de~$X$ sur des points distincts de~$Y$
formant une partie discr\`ete.
\end{enumerate}
\end{remarque}

\section[Th\'eor\`emes de comparaison cohomologique]{Th\'eor\`emes de comparaison cohomologique et th\'eor\`emes d'existence}
\label{XII.4}

L'objet de ce num\'ero est de red\'emontrer les r\'esultats de
\cite[\no 2 th.\, 5 et th.\, 6]{XII.3}; ces derniers g\'en\'eralisent au cas
d'un sch\'ema propre les th\'eor\`emes \'etablis dans \cite[\no 12]{XII.10} lorsque $X$ est projectif, et les \'etendent au cas
relatif. Des r\'esultats plus g\'en\'eraux, concernant les
sch\'emas relatifs propres sur un espace analytique, sont
prouv\'es dans \cite[ch.\, VIII \no 3]{XII.7}.

Rappelons que la cohomologie de {\v{C}}ech utilis\'ee dans \cite[\no 12]{XII.10} co\"incide avec la cohomologie usuelle dans le cas
alg\'ebrique comme dans le cas analytique (EGA~III 1.4.1. et \cite[II 5.10]{XII.5}).
\subsection{}
\label{XII.4.1}
Soient $f \colon X \to Y $ un morphisme de~$\CC$-sch\'emas
localement de type fini et consid\'erons le diagramme commutatif
$$
\xymatrix{{X^\an}\ar[r]^\varphi\ar[d]_{f^\an}&{X}\ar[d]^f\\{Y^\an}\ar[r]^\psi&{Y \quoi.}}
$$
Si $F$ est un $\cal{O}_X$-Module, on a, pour tout entier $p \geq
0$, des morphismes
$$
\R^p f_* F \lto{i} \R^p f_* (\varphi_* F^\an )
\lto{j} \R^p (f\cdot\varphi)_* F^\an \lto{k} \psi_*
(\R^p f_*^\an F^\an ) \quoi,
$$
o\`u $i$ se d\'eduit du morphisme canonique $F \to \varphi_*
F^\an $, et $j, k$ sont des \og edge-homomorphismes\fg de suites
spectrales de Leray. Au compos\'e $k.j.i$ est associ\'e un
morphisme canonique
\begin{equation*}
\label{eq:XII.4.1.1}
\tag{4.1.1} {\theta_p \colon (\R^p f_* F)^\an \to \R^p f_*^\an
(F^\an )}
\end{equation*}

\begin{theoreme}
\label{XII.4.2}
Soient
$f \colon X \to Y $ un morphisme propre de~$\CC$-sch\'emas
localement de type fini, $F$ un $\cal{O}_X$-Module
coh\'erent. Alors, pour tout entier $p \geq 0$ le morphisme
\eqref{eq:XII.4.1.1}
$$
\theta_p \colon (\R^p f_* F)^\an \to \R^p f_*^\an (F^\an)
$$
est un isomorphisme.
\end{theoreme}

1) \emph{Cas où $f$ est projectif}. La d\'emonstration est
analogue \`a celle de \cite[\no 13]{XII.10}. Rappelons-la bri\`evement. On
se ram\`ene au cas o\`u $X$ est un espace projectif type $\PP^r_Y
$ au-dessus de~$Y$. Soit $\cal{Y} = Y^\an$, $\cal{P} = \PP^r_\cal{Y}$;
on prouve d'abord que l'on a
$$
f_*^\an \cal{O}_\cal{P} = \cal{O}_\cal{Y} \quoi, \qquad \R^p f_*^\an
(\cal{O}_\cal{P}) = 0 \qquad \text{pour } p>0
$$
Pour v\'erifier les relations pr\'ec\'edentes, on peut en effet
se ramener au cas o\`u $\cal{Y}$ est une boule $\cal{B}$ d'un espace
affine $\cal{E}^n$. On consid\`ere le \og recouvrement standard\fg $\{
\cal{U}_i \}$ de~$\cal{P}$ par $r+1$ ouverts isomorphes \`a $\cal{B}
\times \cal{E}^r$. Comme ces ouverts sont de Stein, on~a, pour tout
entier $p \geq 0$, des isomorphismes
$$
\H^p (\{ \cal{U}_i \}, \cal{O}_\cal{P}) \isomto \H^p (\cal{P},
\cal{O}_\cal{P})
$$
On peut alors exprimer les sections du faisceau structural
$\cal{O}_\cal{P}$ sur les ouverts $\cal{U}_i $ et sur leurs
intersections en termes de s\'eries de Laurent; un calcul facile
prouve que l'on~a
$$
\H^0 (\cal{P}, \cal{O}_\cal{P} ) \isomto \H^0 (\cal{Y},
\cal{O}_\cal{Y} ) \quoi, \qquad \H^p (\cal{P}, \cal{O}_\cal{P} ) = 0
\qquad \text{pour } p>0
$$
La d\'emonstration s'ach\`eve alors en recopiant \cite[\no 12
lemme~5]{XII.10}, les groupes de cohomologie \'etant remplac\'es par les
faisceaux de cohomologie.

2)
\emph{Cas où $f$ est propre}. On utilise EGA~III 3.1.2 pour se
ramener au cas projectif. Soit $\cal{K}$ la cat\'egorie des
$\cal{O}_X$-Modules coh\'erents tels que $\theta_p$ soit un
isomorphisme pour tout $p \geq 0$. Il suffit de prouver que, pour
toute suite exacte $0 \to F' \to F \to F'' \to 0$
dont deux termes sont dans $\cal{K}$, il en est de m\^eme du
troisi\`eme, qu'un facteur direct d'un objet de~$\cal{K}$ est dans
$\cal{K}$, et que, pour tout point $x$ de~$X$, on peut trouver un
objet $F$ de~$\cal{K}$ tel que l'on ait $F_x \neq 0$.

Le premi\`ere condition r\'esulte par application du lemme des
cinq du diagramme commutatif suivant dont les lignes sont exactes:
$$
\xymatrix@=.7cm{ \ar[r] &{(\R^p f_* F')^\an}\ar[d]\ar[r]&{(\R^p f_*
F)^\an}\ar[d]\ar[r]&{(\R^p f_* F'')^\an}\ar[d]\ar[r]&{(\R^{p+1} f_* F' )^\an}\ar[r]\ar[d]&{\,}\\ \ar[r] &{\R^p
f_*^\an F'^\an}\ar[r]&{\R^p f_*^\an F^\an}\ar[r]&{\R^p
f_*^\an F''^\an}\ar[r]&{\R^{p+1} f_*^\an F'^\an}\ar[r]&{\quoi,}}
$$
et on v\'erifie de fa\c{c}on analogue la deuxi\`eme condition.

Pour v\'erifier la troisi\`eme condition, on peut se borner au cas
o\`u $X$ est un sch\'ema irr\'eductible de point
g\'en\'erique $x$. On pouvait supposer $Y$ noeth\'erien d\`es
le d\'ebut. D'apr\`es le lemme de Chow (EGA~II 5.6.1), on peut
trouver un $Y$-sch\'ema projectif~$X'$ et un morphisme projectif
surjectif $g \colon X' \to X$. D'autre part il existe un entier~$n$
tel que l'on ait $\R^p g_* (\cal{O}_{X'} (n)) = 0 $ pour tout $p>0$ et
que le morphisme canonique $g^* g_* (\cal{O}_{X'} (n)) \to
\cal{O}_{X'} (n)$ soit surjectif (EGA~III 2.2.1). Si l'on pose
$F = g_* (\cal{O}_{X'} (n)) $, le faisceau $F$ r\'epond
\`a la question. En effet on a $F_x \neq 0$; de plus la suite
spectrale de Leray
$$
\R^p f_* (\R^q g_* (\cal{O}_{X'} (n))) \To \R^{p+q}
(f\cdot g)_* (\cal{O}_{X'} (n))
$$
\'etant d\'eg\'en\'er\'ee, on a un isomorphisme
$$
\R^p f_* F \isomto \R^p (f\cdot g)_* (\cal{O}_{X'} (n))
$$
Comme dans le cas alg\'ebrique on a un isomorphisme canonique
$$\R^p f_*^\an F^\an \isomto
\R^p (f\cdot g)_*^\an (\cal{O}_{X'}(n)^\an ),$$
et le diagramme
$$
\xymatrix@C=1cm{{(\R^p f_* F)^\an}\ar[d]_{\theta_p}\ar[r]^-{\sim}&{(\R^p (f\cdot g)_* (\cal{O}_{X'} (n)))^\an}\ar[d]_{\psi_p}\\{\R^p f_*^\an F^\an}\ar[r]^-{\sim}&{\R^p (f\cdot g)_*^\an
(\cal{O}_{X'}(n)^\an)}}
$$
est commutatif. D'apr\`es~1) $\psi_p $ est un isomorphisme; il en
est donc de m\^eme de~$\theta_p$, ce qui ach\`eve la
d\'emonstration.

\begin{corollaire}
\label{XII.4.3}
Soient $X$ un $\CC$-sch\'ema propre, $F$ un $\cal{O}_X$-Module
coh\'erent. Alors, pour tout entier $p\geq 0$, le morphisme
canonique
$$
\H^p(X,F)\to \H^p(X^{\an},F^{\an})
$$
est un isomorphisme.
\end{corollaire}

\begin{theoreme}
\label{XII.4.4}
Soit $X$ un $\CC$-sch\'ema propre. Le foncteur qui, \`a tout
$\cal{O}_X$-Module coh\'erent $F$, associe son image inverse
$F^{\an}$ sur~$X^{\an}$ est une \'equivalence de cat\'egories.
\end{theoreme}

1) \emph{Le foncteur est pleinement fidèle}. Soient en effet $F$
et $G$ deux $\cal{O}_X$-Modules coh\'erents. Le morphisme canonique
$$
\Hom_{\cal{O}_X}(F,G)\to \Hom_{\cal{O}_{X^{\an}}}(F^{\an},G^{\an})
$$
s'identifie au morphisme canonique
$$
\H^0(X,\SheafHom_{\cal{O}_X}(F,G))\to
\H^0(X^{\an},\SheafHom_{\cal{O}_X}(F,G))
$$
(EGA~$0_{\textup{I}}$~6.7.6). Comme $\SheafHom_{\cal{O}_X}(F,G)$ est
coh\'erent, il r\'esulte de~\ref{XII.4.3} que ce morphisme est
bijectif.

2) \emph{Le foncteur est essentiellement surjectif}. Lorsque $X$ est
projectif l'assertion r\'esulte de \cite[\no 12~th.\, 3]{XII.10}. Le cas
g\'en\'eral se ram\`ene au pr\'ec\'edent en utilisant le
lemme de Chow (EGA~II~5.6.1). Soient en effet $X'$ un $\CC$-sch\'ema
projectif, $f\colon X'\to X$ un morphisme projectif surjectif, $U$ un
ouvert dense de~$X$ tel que $f$ induise un isomorphisme
$f^{-1}(U)\simeq U$. On raisonne par r\'ecurrence noeth\'erienne
sur~$X$; on peut donc supposer que, pour tout faisceau coh\'erent
$\cal{G}$ sur~$X^{\an}$ tel que l'on puisse trouver une partie
ferm\'ee $Y$ de~$X$ distincte de~$X$, satisfaisant \`a la relation
$Y^{\an}\supset \Supp \cal{G}$, il existe un faisceau coh\'erent $G$
sur~$X$ tel que l'on ait un isomorphisme $G^{\an}\simeq \cal{G}$.

Soit $\cal{F}$ un faisceau de modules coh\'erent sur~$\cal{O}_{X^{\an}}$, $\cal{K}$ et $\cal{L}$ les faisceaux
coh\'erents d\'efinis par la condition que la suite
$$
0\to\cal{K}\to\cal{F}\to f_*^{\an}f^{\an *}\cal{F}\to\cal{L}\to 0
$$
soit exacte. Comme $X'$ est projectif, il existe un
$\cal{O}_{X'}$-Module coh\'erent $F'$ tel que l'on ait
$F'^{\an}\simeq {f^{\an}}^*\cal{F}$; on d\'eduit alors
de~\ref{XII.4.2} que l'on a un isomorphisme $(f_*F')^{\an}\simeq
f_*^{\an}{f^{\an}}^*\cal{F}$. Comme $\cal{K}|U^{\an}$ et
$\cal{L}|U^{\an}$ sont nuls, il existe des $\cal{O}_X$-Modules
coh\'erents $K$ et $L$ tel que l'on ait des isomorphismes
$K^{\an}\simeq \cal{K}$, $L^{\an}\simeq \cal{L}$. D'apr\`es~1) le
morphisme $f_*^{\an}{f^{\an}}^*\cal{F}\to\cal{L}$ provient d'un unique
morphisme $f_*F'\to L$; soit $I=\Ker(f_*F'\to L)$. Le faisceau
$\cal{F}$ est alors extension de~$I^{\an}$ par~$K^{\an}$, et il suffit
de voir que cette extension provient
par image inverse d'une extension de~$I$ par~$K$. Il suffit donc de
prouver que le morphisme canonique
\begin{equation*}
\label{eq:XII.4.*}
\tag{$*$} \Ext_{\cal{O}_X}^q(I,K)^{\an}\isomto
\Ext_{\cal{O}_{X^{\an}}}^q(I^{\an},K^{\an})\quad q\neq 1
\end{equation*}
est bijectif. Or on a des isomorphismes
$\mathbf{Ext}_{\cal{O}_X}^q(I,K)^{\an}\isomto
\mathbf{Ext}_{\cal{O}_{X^{\an}}}^q(I^{\an},K^{\an})$ pour tout entier
$q\geq 0$ (EGA~$0_{\textup{III}}$~12.3.5), et un morphisme de suites
spectrales
$$
\xymatrix@R=.7cm@C=.5cm{{\H^p(X,\mathbf{Ext}_{\cal{O}_X}^q(I,K))}\ar@{=>}[r]\ar[d]&{\mathrm{Ext}_{\cal{O}_X}^{p+q}(I,K)}\ar[d]\\{\H^p(X^{\an},\mathbf{Ext}_{\cal{O}_{X^{\an}}}^q(I^{\an},K^{\an}))}\ar@{=>}[r]&{\mathrm{Ext}_{\cal{O}_{X^{\an}}}^{p+q}(I^{\an},K^{\an}).}}
$$
Ce morphisme est un isomorphisme car, d'apr\`es~\ref{XII.4.3}, il en
est ainsi sur les termes $E^{pq}_2$, et ceci d\'emontre la
bijectivit\'e de~$(*)$.

\begin{corollaire}
\label{XII.4.5}
Le foncteur qui \`a tout $\CC$-sch\'ema propre $X$ associe
$X^{\an}$ est pleinement fid\`ele.
\end{corollaire}

On doit montrer que, si $X$ et $Y$ sont deux $\CC$-sch\'emas
propres, l'application canonique
$$
\Hom_{\CC}(X,Y)\to\Hom(X^{\an},Y^{\an})
$$
est bijective. Or se donner un morphisme de~$X$ dans~$Y$ (\resp de
$X^{\an}$ dans $Y^{\an}$) \'equivaut \`a se donner son graphe,
\ie un sous-sch\'ema ferm\'e $Z$ de~$X\times Y$ (\resp un
sous-espace analytique ferm\'e $\othercal{Z}$ de~$X^{\an}\times
Y^{\an}$), tel que la restriction de la premi\`ere projection
$X\times Y\to X$ \`a~$Z$ (\resp de~$X^{\an}\times Y^{\an}\to
X^{\an}$ \`a~$\othercal{Z}$) soit
un isomorphisme. Comme la donn\'ee d'un sous-sch\'ema ferm\'e
de~$X\times Y$ (\resp d'un sous-espace analytique ferm\'e
de~$X^{\an}\times Y^{\an}$) \'equivaut \`a celle d'un faisceau
coh\'erent d'id\'eaux sur~$\cal{O}_{X\times Y}$ (\resp sur~$\cal{O}_{X^{\an}\times Y^{\an}}$), le corollaire r\'esulte
de~\ref{XII.4.4}.

\begin{corollaire}
\label{XII.4.6}
Soit $X$ un $\CC$-sch\'ema propre. Le foncteur qui, \`a tout
sch\'ema fini (\resp \'etale fini) $X'$ au-dessus de~$X$, associe
$X'^{\an}$ est une \'equivalence de la cat\'egorie des sch\'emas
finis (\resp \'etales finis) au-dessus de~$X$ dans la cat\'egorie
des espaces analytiques finis (\resp \'etales finis) au-dessus
de~$X^{\an}$.
\end{corollaire}

En effet se donner un morphisme fini $X'\to X$ (\resp $X'^{\an}\to
X^{\an}$) \'equivaut \`a se donner un faisceau coh\'erent
d'alg\`ebres sur~$\cal{O}_X$ (\resp sur~$\cal{O}_{X^{\an}}$)
\cite[\no 19~\S 5~th.\, 2]{XII.4}. Le corollaire r\'esulte donc de~\ref{XII.4.4}
dans le cas non resp\'e, et le cas resp\'e s'en d\'eduit compte
tenu de~\ref{XII.3.1}~(iii).

\section{Th\'eor\`emes de comparaison des rev\^etements \'etales}
\label{XII.5}
\setcounter{subsection}{-1}
\subsection{}
\label{XII.5.0}
Pr\'ecisons la notion de rev\^etement fini d'un espace analytique.
\index{revetement fini d'un espace analytique@rev\^etement fini d'un espace analytique|hyperpage}Si $\othercal{X}$ est un espace analytique, on dit qu'un espace analytique
$\othercal{X}'$ fini au-dessus de~$\othercal{X}$ est un rev\^etement fini de
$\othercal{X}$ si toute composante irr\'eductible de~$\othercal{X}'$ domine
une composante irr\'eductible de~$\othercal{X}$.


\begin{theoreme}[\og Th\'eor\`eme d'existence de Riemann\fg]
\label{XII.5.1}
\index{Riemann (theoreme d'existence de)@Riemann (th\'eor\`eme d'existence de)|hyperpage}\index{theoreme d'existence de Riemann@ th\'eor\`eme d'existence de Riemann|hyperpage}Soient $X$ un $\CC$-sch\'ema localement de type fini, $X^{\an}$
l'espace analytique associ\'e \`a $X$. Le foncteur $\Psi$ qui,
\`a tout rev\^etement \'etale fini $X'$ de~$X$, associe
$X^{\prime\an}$ est une \'equivalence
de la cat\'egorie des rev\^etements \'etales finis de~$X$ dans
la cat\'egorie des rev\^etements \'etales finis de~$X^{\an}$.
\end{theoreme}

1) \emph{Le foncteur $\Psi$ est pleinement fidèle.} Soient $X'$
et $X''$ deux rev\^etements \'etales finis de~$X$, et prouvons que
l'application canonique
\begin{equation*}
\label{eq:XII.5.*}
\tag{$*$} { \Hom_{X}(X',X'') \to
\Hom_{X^{\an}}(X^{\prime\an},X^{\prime\prime\an}) }
\end{equation*}
est bijective. On peut supposer $X'$ connexe. Se donner un
$X$-morphisme de~$X'$ dans $X''$ \'equivaut \`a se donner une
composante connexe $X_{i}$ de~$X'\times_{X}X''$ telle que le morphisme
$X_{i}\to X'$ induit par la premi\`ere projection soit un
isomorphisme. Comme les composantes connexes de~$X'\times_{X}X''$
correspondent bijectivement aux composantes connexes de
$X^{\prime\an}\times_{X^{\an}}X^{\prime\prime\an}$~\eqref{XII.2.6} et qu'un
morphisme $X_{i}\to X'$ est un isomorphisme si et seulement si il en
est ainsi de~$X_{i}^{\an}\to X^{\prime\an}$, ceci d\'emontre la
bijectivit\'e de~\eqref{eq:XII.5.*}.

2) \emph{Le foncteur $\Psi$ est essentiellement surjectif.} Soit
$\othercal{X}'$ un rev\^etement \'etale fini de~$X^{\an}$ et prouvons
qu'il existe un rev\^etement \'etale $X'$ de~$X$ tel que l'on ait
un isomorphisme $X^{\prime\an}\isomto\othercal{X}'$. Compte tenu de~1) la
question est locale sur~$X$, et on peut donc supposer $X$ affine.

a) R\'eduction au cas o\`u $X$ est normal. On peut supposer $X$
r\'eduit. Supposons en effet le th\'eor\`eme d\'emontr\'e
pour $X_{\red}$. Le foncteur qui, \`a un rev\^etement \'etale
fini $X'$ de~$X$ fait correspondre le rev\^etement \'etale fini
$X^{\prime\an}_{\red}$ de~$X_{\red}^{\an}$ est alors une \'equivalence.
Comme il s'obtient en composant $\Psi$ avec le foncteur $\Theta$ qui,
\`a un rev\^etement \'etale fini de~$X^{\an}$ associe son image
inverse sur~$X_{\red}^{\an}$, et que $\Theta$ est pleinement
fid\`ele, ceci montre que $\Psi$ est une \'equivalence de
cat\'egories.

On
peut supposer $X$ normal. Soit en effet $\widetilde{X}$ le
normalis\'e de~$X$, $p\colon\widetilde{X}\to X$ le morphisme
canonique. Comme $p$ est fini, $p$ est un morphisme de descente
effective pour la cat\'egorie des rev\^etements
\'etales~(IX~\ref{IX.4.7}). Le th\'eor\`eme \'etant
suppos\'e d\'emontr\'e pour $\widetilde{X}$, si l'on pose
$\widetilde{\othercal{X}'}=\othercal{X}'\times_{X^{\an}}\widetilde{X}^{\an}$,
il existe un rev\^etement \'etale $\widetilde{X}'$ de
$\widetilde{X}$ et un isomorphisme $\widetilde{X}'{}^{\an}\simeq
\widetilde{\othercal{X}'}$. Il r\'esulte alors de~1) que la donn\'ee
de descente naturelle que l'on a sur~$\widetilde{\othercal{X}'}$ se
rel\`eve en une donn\'ee de descente sur~$\widetilde{X}'$
relativement \`a $\widetilde{X}\to X$; ceci prouve l'existence d'un
rev\^etement \'etale $X'$ de~$X$ tel que l'on ait un isomorphisme
$i\colon X^{\prime\an}\times_{X^{\an}}\widetilde{X}^{\an} \simeq
\widetilde{\othercal{X}'}$, dont les images inverses par les deux
projections de
$\widetilde{X}^{\an}\times_{X^{\an}}\widetilde{X}^{\an}$ soient les
m\^emes. D'apr\`es~IX~\ref{IX.3.2}, dont la d\'emonstration est
valable dans le cas analytique, le morphisme $\widetilde{X}^{\an}\to
X^{\an}$ est un morphisme de descente pour la cat\'egorie des
rev\^etements \'etales, et par suite $i$ provient d'un
isomorphisme $X^{\prime\an}\simeq
\othercal{X}'$.

b) R\'eduction au cas o\`u $X$ est r\'egulier. Soient $U$
l'ouvert des points r\'eguliers de~$X$, $i\colon U\to X$,
$i^{\an}\colon U^{\an}\to X^{\an}$ les morphismes canoniques; comme
$X$ est normal, on a $\codim(X-U,X)\geq 2$. Supposons qu'il existe un
rev\^etement \'etale $U'$ de~$U$ tel que l'on ait
$U^{\prime\an}\simeq \othercal{X}'|U^{\an}$ et montrons qu'alors $U'$ se
prolonge en un rev\^etement \'etale $X'$ de~$X$ tel que l'on ait
$X^{\prime\an}\simeq \othercal{X}'$. Il suffit de voir que $U'$ se prolonge
en un rev\^etement \'etale $X'$ de~$X$; en effet on aura alors un
isomorphisme $X^{\prime\an}|U^{\an}\simeq \othercal{X}'|U^{\an}$; mais, si
$\cal{F}$ et $\cal{G}$ sont les faisceaux coh\'erents d'alg\`ebres
sur~$\cal{O}_{X^{\an}}$ d\'efinissant respectivement $\othercal{X}'$ et
$X^{\prime\an}$, le fait que $X$ soit normal et que l'on ait
$\codim(X-U,X)\geq 2$ entra\^ine que les morphismes canoniques
$$
\cal{F}\to i_*^{\an}(\cal{F}|U^{\an}) \qquad
\cal{G}\to i_*^{\an}(\cal{G}|U^{\an})
$$
sont
des isomorphismes \cite[\no 3~prop.\, 4]{XII.11}. Il en r\'esulte que $\cal{F}$
et $\cal{G}$ donc aussi $X^{\prime\an}$ et $\othercal{X}'$ sont isomorphes.

Soit $\varphi\colon X^{\an}\to X$ le morphisme canonique. Comme le
probl\`eme de prolonger $U'$ \`a $X$ est local sur~$X$, il suffit
de prouver que, pour tout point $y$ de~$X^{\an}- U^{\an}$, le rev\^etement \'etale
$U'_{\varphi(y)}=U'\times_{X}\Spec\cal{O}_{X,\varphi(y)}$ de
$U_{\varphi(y)}=U\times_{X}\Spec\cal{O}_{X,\varphi(y)}$ se prolonge
\`a $\Spec\cal{O}_{X,\varphi(y)}$. Soit $H$ la
$\cal{O}_{U}$-Alg\`ebre coh\'erente d\'efinissant $U'$. Le
morphisme canonique
$$
\alpha\colon (i_*H)^{\an} \to i_*^{\an}(H^{\an})
=\cal{F}
$$
d\'efinit un morphisme de faisceaux de modules sur~$\Spec\cal{O}_{X^{\an},y}$:
$$
\alpha_{y}\colon (i_*H)^{\an}_{y} \to \cal{F}_{y}\quoi,
$$
dont la restriction \`a $U_{y}=U_{\varphi(y)}
\times_{\Spec\cal{O}_{X,\varphi(y)}}\Spec\cal{O}_{X^{\an},y}$ est un
isomorphisme. Mais ceci prouve que $H|U_{y}$ est trivial, donc que
$U'_{\varphi(y)}$ se prolonge \`a $\Spec\cal{O}_{X,\varphi(y)}$.

c) Cas o\`u $X$ est affine r\'egulier. Soit
$$
j\colon X\to P
$$
une compactification de~$X$, o\`u $P$ est un $\CC$-sch\'ema
projectif et $j$ une immersion ouverte dominante. Gr\^ace au
th\'eor\`eme de r\'esolution des singularit\'es~\cite{XII.8}, on peut
trouver un sch\'ema r\'egulier $R$, un morphisme projectif
$r\colon R\to P$, tel que $r$ induise un isomorphisme $r^{-1}(X)\simeq
X$ et que $r^{-1}(X)$ soit le compl\'ementaire dans $R$ d'un
diviseur \`a croisements normaux. Soit
$$
k\colon X\to R
$$
l'immersion canonique. On va montrer qu'il existe un rev\^etement
fini normal~\eqref{XII.5.0} $\othercal{R}'$ de~$R^{\an}$ qui prolonge le
rev\^etement \'etale $X^{\prime\an}$. D'apr\`es la
proposition~\ref{XII.5.3} ci-dessous, un tel rev\^etement est
unique; le probl\`eme de prolonger $X^{\prime\an}$ est donc local sur~$R^{\an}$ au voisinage de~$R^{\an}- X^{\an}$. Or chaque
point de~$R^{\an}- X^{\an}$ a un voisinage ouvert
$\othercal{V}$ isomorphe \`a une boule d'un espace affine
$\othercal{E}^{n}$, tel que $\othercal{V}- \othercal{V}\cap X^{\an}$
soit d\'efini par l'annulation des $p$ premi\`eres fonctions
coordonn\'ees $z_{1},\dots,z_{p}$, avec $0\leq p\leq n$. Le groupe
fondamental de~$\othercal{U}=\othercal{V}\cap X^{\an}$ est isomorphe \`a
$\ZZ^{p}$, et tout rev\^etement \'etale de~$\othercal{U}$ est quotient
d'un rev\^etement de la forme
$$
\othercal{U}''= \othercal{U}[T_{1},\dots,T_{p}]/
(T_{1}^{n_{1}}-z_{1},\dots,T_{p}^{n_{p}}-z_{p})\quoi,
$$
o\`u les $n_{i}$ sont des entiers $>0$, par un sous-groupe $H$ du
groupe de Galois $\ZZ/n_{1}\ZZ\times\cdots\times\ZZ/n_{p}\ZZ$ de
$\othercal{U}''$. Or $\othercal{U}''$ se prolonge en le rev\^etement
r\'egulier
$$
\othercal{V}''= \othercal{V}[T_{1},\dots,T_{p}]/
(T_{1}^{n_{1}}-z_{1},\dots,T_{p}^{n_{p}}-z_{p})\quoi,
$$
de~$\othercal{V}$ sur lequel $H$ op\`ere, et le quotient de~$\othercal{V}''$
par $H$ est le prolongement cherch\'e.

La d\'emonstration s'ach\`eve alors gr\^ace \`a~\ref{XII.4.6}.
Le rev\^etement $\othercal{R}'$ provient d'un rev\^etement fini $R'$
de~$R$; la restriction de~$R'$ \`a $X$ est un rev\^etement $X'$ de
$X$ tel que l'on ait $X^{\prime\an}\simeq \othercal{X}'$, et
d'apr\`es~\ref{XII.3.1}(iii) $X'$ est un rev\^etement \'etale
de~$X$.

\begin{corollaire}
\label{XII.5.2}
Soient
$X$ un $\CC$-sch\'ema localement de type fini connexe,
$\varphi\colon X^{\an}{\to} X$ le morphisme canonique, $x$ un point de
$X^{\an}$. Soit $\pi_{1}(X^{\an},x)$ le groupe fondamental de
l'espace topologique $X^{\an}$ au point $x$, $\pi_{1}(X,\varphi(x))$
le groupe fondamental du sch\'ema $X$ au point $\varphi(x)$
\textup{(V~\ref{V.7})}. Alors $\pi_{1}(X,\varphi(x))$ est canoniquement
isomorphe au compl\'et\'e de~$\pi_{1}(X^{\an},x)$, pour la
topologie des sous-groupes d'indice fini.
\end{corollaire}

Soit en effet $\cal{C}$ la cat\'egorie des rev\^etements
\'etales finis de~$X^{\an}$, $F$ le foncteur de~$\cal{C}$ dans
$\Ens$ qui, \`a tout rev\^etement \'etale fini $\othercal{X}'$ de
$X^{\an}$ associe l'ensemble des points de~$\othercal{X}'$ au-dessus de
$x$, et soit $\widehat{\pi}_{1}(X^{\an},x)$ le groupe profini
associ\'e \`a $\cal{C}$ et $F$ comme il est dit dans~V~\ref{V.4}.
Comme tout rev\^etement \'etale fini de~$X^{\an}$ est quotient du
rev\^etement universel par un sous-groupe d'indice fini,
$\widehat{\pi}_{1}(X^{\an},x)$ n'est autre que le compl\'et\'e de
$\pi_{1}(X^{\an},x)$ pour la topologie des sous-groupes d'indice fini.
Le corollaire r\'esulte donc de~\ref{XII.5.1} et~V~\ref{V.6.10}.

\begin{proposition}
\label{XII.5.3}
Soient $\othercal{X}$ un espace analytique normal, $\othercal{Y}$ un
sous-ensemble analytique ferm\'e tel que
$\othercal{U}=\othercal{X}-\othercal{Y}$ soit dense dans $\othercal{X}$.
Alors le foncteur qui, \`a tout rev\^etement normal
fini~\eqref{XII.5.0} $\othercal{X}'$ de~$\othercal{X}$ associe sa restriction
\`a $\othercal{U}$ est pleinement fid\`ele.
\end{proposition}

Soient $\othercal{X}'$ and $\othercal{X}''$ deux rev\^etements finis normaux
de~$\othercal{X}$. On doit montrer que l'application canonique
$$
\Hom_{\othercal{X}}(\othercal{X}',\othercal{X}'') \to
\Hom_{\othercal{U}}(\othercal{X}'|\othercal{U},\othercal{X}''|\othercal{U})
$$
est bijective. Soient $u$, $v$ deux $\othercal{X}$-morphismes de
$\othercal{X}'$ dans $\othercal{X}''$ dont les restrictions \`a $\othercal{U}$
sont les m\^emes et prouvons que $u=v$. Les morphismes $u$ et
$v$ co\"incident sur l'ouvert dense
$\othercal{U}\times_{\othercal{X}}\othercal{X}'$, donc sur les espaces
topologiques sous-jacents. D'apr\`es \hbox{\cite[\no 19~\S 4~cor.\, 5]{XII.4}} ceci
prouve que l'on a $u=v$.

Soit maintenant $u$ un $\othercal{U}$-morphisme de~$\othercal{X}'|\othercal{U}$
dans $\othercal{X}''|\othercal{U}$ et montrons qu'il se prolonge \`a
$\othercal{X}'$ tout entier. On peut supposer $\othercal{X}'$ r\'egulier.
En effet, $\othercal{X}'$ \'etant normal, on peut trouver un ouvert
$\othercal{V}$ de~$\othercal{X}$ dont le compl\'ementaire soit une partie
analytique de codimension $\geq 2$, tel que
$\othercal{X}'\times_{\othercal{X}}\othercal{V}=\othercal{V}'$ soit r\'egulier.
Soit $\othercal{V}''=\othercal{X}''\times_{\othercal{X}}\othercal{V}$ et supposons la
proposition d\'emontr\'ee pour $\othercal{V}$. On consid\`ere le
diagramme commutatif
$$
\xymatrix{{\othercal{V}'}\ar[dd]_{g'}\ar[rd]\ar[rrr]^{i'}&& &{\othercal{X}'}\ar[dd]^{f'}\\ &{\othercal{V}''}\ar[ld]_{g''\mkern-10mu}\ar[rrr]^-{i''~~}& &
&{\othercal{X}''}\ar[ld]_{f''\mkern-12mu}\\{\othercal{V}}\ar[rrr]^{i}&& &{\othercal{X}}\\
}
$$
\`A $u$ est associ\'e un morphisme de
$\cal{O}_{\othercal{V}}$-Alg\`ebres $g''_*\cal{O}_{\othercal{V}''}\to
g'_*\cal{O}_{\othercal{V}'}$, d'o\`u l'on d\'eduit un morphisme
$$
i_*g''_*\cal{O}_{\othercal{V}''} \to
i_*g'_*\cal{O}_{\othercal{V}'}\quoi.
$$
Compte tenu des isomorphismes
$i'_*\cal{O}_{\othercal{V}'}\simeq\cal{O}_{\othercal{X}'}$,
$i''_*\cal{O}_{\othercal{V}''}\simeq\cal{O}_{\othercal{X}''}$
\cite[\no 3~prop.\, 4]{XII.11} on en d\'eduit un morphisme de
$\cal{O}_{\othercal{X}}$-Alg\`ebres
$$
f''_*\cal{O}_{\othercal{X}''} \to
f'_*\cal{O}_{\othercal{X}'}\quoi,
$$
d'o\`u le morphisme $\othercal{X}'\to\othercal{X}''$ cherch\'e.

On
suppose d\'esormais $\othercal{X}'$ r\'egulier. Soient
$\othercal{U}'=\othercal{U}\times_{\othercal{X}}\othercal{X}'$,
$\othercal{Y}'=\othercal{X}'-\othercal{U}'$. On consid\`ere $\othercal{Y}'$ comme
sous-espace analytique r\'eduit de~$\othercal{X}'$; si $\othercal{Y}'_{1}$
est le ferm\'e singulier de~$\othercal{Y}'$, on a $\dim \othercal{Y}'_{1} <
\dim \othercal{Y}'$~\cite[\no 20~D.\, Th.\, 3]{XII.4}. On voit donc par r\'ecurrence
sur la dimension de~$\othercal{Y}'$ que l'on peut supposer $\othercal{Y}'$
lisse. Comme il suffit de prolonger $u$ \`a un voisinage ouvert de
chaque point de~$\othercal{Y}'$, on peut supposer par le th\'eor\`eme
des fonctions implicites que $\othercal{X}'$ est une boule d'un espace
affine $\othercal{E}^{n}$ et $\othercal{Y}'$ le ferm\'e d\'efini par
l'annulation des $p$ premi\`eres fonctions coordonn\'ees
$z_{1},\dots,z_{p}$, avec $0\leq p\leq n$.

On associe \`a $u$ une section $s$ de
$p\colon\othercal{X}'\times_{\othercal{X}}\othercal{X}''\to\othercal{X}'$ au-dessus de
$\othercal{U}'$; quitte \`a restreindre $\othercal{X}'$, on peut supposer
$p_*(\cal{O}_{\othercal{X}'\times_{\othercal{X}}\othercal{X}''})$
engendr\'e par des \'el\'ements $x_{1},\dots,x_{q}$ de
$\Gamma(\othercal{X}',
p_*\cal{O}_{\othercal{X}'\times_{\othercal{X}}\othercal{X}''})$; soient
$u_{1},\dots,u_{q}\in\Gamma(\othercal{U}',\cal{O}_{\othercal{X}'})$ les images
par $s$ de~${x_{1}}|\othercal{U}',\dots,{x_{q}}|\othercal{U}'$. Dire
que $s$ se prolonge \`a $\othercal{X}'$ revient \`a dire que les
$u_{1},\dots,u_{q}$ se prolongent en des sections de
$\Gamma(\othercal{X}',\cal{O}_{\othercal{X}'})$. Mais, puisque $f$ est fini,
chaque $u_{i}$ est une s\'erie de Laurent en $z_{1},\dots,z_{p}$,
\`a coefficients des s\'eries enti\`eres en
$z_{p{+}1},\dots,z_{n}$, qui satisfont \`a des relations de
d\'ependance int\'egrale. Il en r\'esulte que $u_{i}$ est
born\'e donc est une s\'erie enti\`ere en $z_{1},\dots,z_{n}$,
et par suite se prolonge \`a $\othercal{X}'$.

On peut se demander si le foncteur introduit dans~\ref{XII.5.3} est
une \'equivalence de cat\'egories. On a une r\'eponse \`a
cette question gr\^ace au th\'eor\`eme de \textsc{Grauert}-\textsc{Remmert} \cite{XII.6}
dont nous donnons une d\'emonstration ci-dessous utilisant la
r\'esolution des singularit\'es. On aurait aussi pu utiliser le
th\'eor\`eme de \textsc{Grauert}-\textsc{Remmert} pour d\'emontrer~\ref{XII.5.1};
c'est ce que l'on faisait avant de disposer de~\cite{XII.8}.

\enlargethispage{.5cm}
\begin{theoreme}[Th\'eor\`eme de \textsc{Grauert}-\textsc{Remmert}]
\label{XII.5.4}
Soient
\index{Grauert-Remmert (th\'eor\`eme de)|hyperpage}\index{theoreme de Grauert-Remmert@th\'eor\`eme de Grauert-Remmert|hyperpage}$\othercal{X}$ un espace analytique normal, $\othercal{Y}$ un
sous-ensemble analytique ferm\'e tel que
$\othercal{U}=\othercal{X}-\othercal{Y}$ soit dense dans~$\othercal{X}$. Soit
$\othercal{U}'$ un rev\^etement normal fini de~$\othercal{U}$; on suppose
qu'il existe une partie analytique ferm\'ee rare $\othercal{S}$
de~$\othercal{X}$ telle que la restriction de~$\othercal{U}'$ \`a
$\othercal{U}-\othercal{U}\cap\othercal{S}$ soit \'etale. Alors il existe un
rev\^etement fini normal $\othercal{X}'$ de~$\othercal{X}$ qui prolonge
$\othercal{U}'$, et $\othercal{X}'$ est unique \`a isomorphisme pr\`es.
\end{theoreme}

L'unicit\'e r\'esulte de~\ref{XII.5.3}. Le probl\`eme de
prolonger $\othercal{U}'$ est donc local sur~$\othercal{X}$. On peut supposer
$\othercal{U}$ r\'egulier et $\othercal{U}'$ \'etale sur~$\othercal{U}$. En
effet l'ensemble des points r\'eguliers de~$\othercal{U}$ est un ouvert
$\othercal{V}$ dense dans~$\othercal{X}$ dont le compl\'ementaire est une
partie analytique \hbox{\cite[\no20~D th.\, 2]{XII.4}} et il suffit de remplacer
$\othercal{U}$ par l'ouvert $\othercal{V}-\othercal{V}\cap\othercal{S}$.

Soit $y$ un point de~$\othercal{X}-\othercal{U}$ et montrons que l'on peut
prolonger $\othercal{U}'$ \`a un voisinage de~$y$. Quitte \`a
restreindre $\othercal{X}$ \`a un voisinage ouvert de~$y$, il
r\'esulte du th\'eor\`eme de r\'esolution des singularit\'es
\cite{XII.8} que l'on peut trouver un espace analytique r\'egulier
$\othercal{X}_1$, un morphisme projectif $f\colon\othercal{X}_1\to\othercal{X}$
induisant par restriction \`a~$\othercal{U}$ un isomorphisme
$\othercal{U}_1=f^{-1}(\othercal{U})\simeq\othercal{U}$, tel que $\othercal{U}_1$ soit
le compl\'ementaire dans~$\othercal{X}_1$ d'un diviseur \`a
croisements normaux. Montrons que $\othercal{U}'$ se prolonge en un
rev\^etement fini normal de~$\othercal{X}_1$. Comme la question est
locale sur~$\othercal{X}_1$, on peut supposer que $\othercal{X}_1$ est une
boule d'un espace affine $\othercal{E}^n$ et que $\othercal{X}_1-\othercal{U}_1$
est d\'efini par l'annulation des $p$ premi\`eres fonctions
coordonn\'ees
$z_1,\dots,z_p$,
avec $0\leq p\leq n$. Le rev\^etement \'etale $\othercal{U}'$
de~$\othercal{U}_1$ est quotient d'un rev\^etement de la forme
$$
\othercal{U}_2=\othercal{U}_1[T_1,\dots,T_p]/\left(T_1^{n_1}-z_1,\dots,
T_p^{n_p}-z_p\right)
$$
par un sous-groupe $H$ du groupe de Galois de~$\othercal{U}_2$. Le
rev\^etement $\othercal{U}_2$ se prolonge
en le rev\^etement
$$
\othercal{X}_2=\othercal{X}_1[T_1,\dots,T_p]/\left(T_1^{n_1}-z_1,\dots,
T_p^{n_p}-z_p\right)
$$
de~$\othercal{X}_1$ sur lequel $H$ op\`ere, et $\othercal{X}_2/H$ prolonge
$\othercal{U}'$ \`a~$\othercal{X}_1$.

Notons $\othercal{X}_1'$ le rev\^etement fini normal de~$\othercal{X}_1$ qui
prolonge $\othercal{U}'$, $\othercal{F}_1$ la
$\cal{O}_{\othercal{X}_1}$-Alg\`ebre coh\'erente d\'efinie
par~$\othercal{F}_1$. D'apr\`es le th\'eor\`eme de finitude de
\textsc{Grauert}-\textsc{Remmert} \hbox{\cite[\no15 th.\, 1.1]{XII.4}}, $f_*\othercal{F}_1$ est une
$\cal{O}_{\othercal{X}}$-Alg\`ebre coh\'erente. Il lui correspond donc
un rev\^etement fini $\othercal{X}'$ de~$\othercal{X}$ qui est d'ailleurs
normal puisque $\othercal{X}_1'$ l'est, et $\othercal{X}'$ est le prolongement
de~$\othercal{U}'$ cherch\'e.

\begin{remarque}
\label{XII.5.5}
Dans l'\'enonc\'e~\ref{XII.5.4}, on ne peut supprimer
l'hypoth\`ese sur le lieu des points o\`u le morphisme
$\othercal{U}'\to\othercal{U}$ n'est pas \'etale. Soit par exemple
$\othercal{X}$ le disque unit\'e du plan complexe, $\othercal{U}$ le
compl\'ementaire de l'origine dans~$\othercal{X}$,
$\othercal{U}'=\othercal{U}[T]/(T^2-\sin 1/z)$, o\`u~$z$ est la fonction
coordonn\'ee sur~$\othercal{X}$. Alors $\othercal{U}'$ est un rev\^etement
fini normal de~$\othercal{U}$ qui ne se prolonge pas
\`a~$\othercal{X}$. Supposons en effet que $\othercal{U}'$ se prolonge en un
rev\^etement fini $\othercal{X}'$ de~$\othercal{X}$; le lieu des points
de~$\othercal{X}$ o\`u le morphisme $\othercal{X}'\to\othercal{X}$ n'est pas
\'etale est alors un ferm\'e analytique qui contient tous les
points $z$ tels que l'on ait $\sin 1/z=0$, ce qui est absurde.

On peut cependant supprimer l'hypoth\`ese sur le lieu singulier du
morphisme $\othercal{U}'\to\othercal{U}$ lorsque l'on a
$\codim(\othercal{X}-\othercal{U},\othercal{X})\geq~2$. On peut en effet supposer
$\othercal{U}$ r\'egulier. Le lieu des points de~$\othercal{U}$ o\`u
$\othercal{U}'\to\othercal{U}$ n'est pas \'etale est un diviseur
de~$\othercal{U}$, et il r\'esulte du th\'eor\`eme de \textsc{Remmert}-\textsc{Stein}
\cite[th.\, 3]{XII.9} qu'il est
la trace sur~$\othercal{U}$ d'un diviseur de~$\othercal{X}$. Or, dans ce cas,
si $\othercal{A}$ est une $\cal{O}_{\othercal{U}}$-Alg\`ebre coh\'erente
telle que $\othercal{U}'=\Spec \an(\othercal{A})$, si $i:\othercal{U}\to\othercal{X}$
est le morphisme canonique, il suffit de prendre
$\othercal{X}'=\Spec\an(i_*\othercal{A})$; on sait en effet que $i_*\othercal{A}$
est
coh\'erente
\cite[\no 1 th.\, 1]{XII.11}.
\end{remarque}

\begin{remarqueMR}
\label{XII.5.6}
Il existe des groupes de pr\'esentation finie $G$, non
triviaux, qui ne poss\`edent pas de sous-groupes d'indice fini,
distincts de~$G$, par exemple le groupe de G.\, Higman (\cf J-P.\, Serre, Arbres et amalgames, Ast\'erisque N$^\circ$46, prop.\, 6, chap.\, I, \S 1).
D\`es lors, si un tel groupe se r\'ealise comme groupe fondamental
topologique d'un sch\'ema $V$ sur~$\CC$, disons lisse et
projectif, l'espace topologique $V_{\mathrm{top}}$ sous-jacent \`a $V$ n'est pas
simplement connexe, mais $V$ est alg\'ebriquement simplement connexe.
\`A l'heure actuelle, on ne conna\^it pas de tels $V$. Signalons toutefois
que D.\, Toledo a construit des sch\'emas lisses et projectifs $V$ sur~$\CC$ dont le groupe fondamental topologique n'est pas s\'epar\'e
pour la topologie des sous-groupes d'indice fini (le morphisme naturel
$\pi_1(V_{\mathrm{top}})\to\pi_1(V_{\alg})$ n'est pas injectif). [D.\, Toledo, Projective varieties with non-residually finite fundamental group, Publ.\ Math.\ IHES \textbf{77} (1993), p.\,103--119].
\end{remarqueMR}


\begin{thebibliography}{10}{XII.6}

\bibitem[1]{XII.1} N.\, \textsc{Bourbaki},
Topologie G\'en\'erale, Hermann, Paris (1960).

\bibitem[2]{XII.2} N.\, \textsc{Bourbaki},
Alg\`ebre Commutative, Hermann, Paris (1961).

\bibitem[3]{XII.3} H.\, \textsc{Cartan}. S\'eminaire E.N.S., Paris (1956-57).

\bibitem[4]{XII.4} H.\, \textsc{Cartan}. S\'eminaire E.N.S., Paris (1960-61).

\bibitem[5]{XII.5} R.\, \textsc{Godement}, Th\'eorie des Faisceaux, Hermann, Paris
(1958).

\bibitem[6]{XII.6} H.\, \textsc{Grauert} et R.\, \textsc{Remmert},
Komplexe R\"aume, Math.\ Ann. \textbf{136} (1958), p.\, 245--318.

\bibitem[7]{XII.7} M.\, \textsc{Hakim}, Sch\'emas relatifs, th\`ese, Paris (1967).

\bibitem[8]{XII.8} H.\, \textsc{Hironaka},
Resolution of singularities of an algebraic variety over a field of characteristic zero, Ann.\ of Math. \textbf{39} (1964), p.\, 109--236.


\bibitem[9]{XII.9} R.\, \textsc{Remmert} et K.\, \textsc{Stein},
Ueber die wesentlichen Singularit\"aten analytischer Mengen,
Math.\ Ann. \textbf{126} (1953), p.\, 263--306.

\bibitem[10]{XII.10}
J-P.\,\textsc{Serre},
G\'eom\'etrie alg\'ebrique et g\'eom\'etrie analytique, Ann.\ Inst.\ Fourier (Grenoble) \textbf{6} (1956), p.\, 1--42.

\bibitem[11]{XII.11}
J-P.\,\textsc{Serre},
Prolongement de faisceaux analytiques
coh\'erents, Ann.\ Inst.\ Fourier (Grenoble) \textbf{16} (1966), p.\, 363--374.


\end{thebibliography}

\chapterspace{-4}
